\section{What Can Be Quantified as Risk?}\label{sec:risk}

The defensible quantitative outputs are conditional growth profiles,
uncertainty intervals, observed banking-stress interactions, and held-out
forecast losses. At an observed stress value $f$, the reported profile is
$\hat\beta_h+\hat\phi_{F,h}f$, with other conditioning states at their
reference values. It should be evaluated within the empirical support and
with its covariance-based uncertainty. A change in this profile is not an
estimated intervention effect.

The earlier twelve-year growth scenarios are withdrawn. The original code
convolved cumulative local-projection coefficients with repeated deviations
in the index level and divided the result by $\sqrt{H}$. Neither overlapping
cumulative coefficients nor this scaling identifies annual growth responses
to a sequence of structural shocks. Differencing cumulative coefficients
alone would not repair the additional problem that the regressor is a
persistent, endogenous level rather than an identified innovation.
A replacement scenario system would require explicit dynamic equations,
shock identification, and validation of the mapping from interventions to
parameters. Those ingredients are not supplied by the current panel.

The composite burden score is also withdrawn from quantitative claims.
Normalizing a channel by its estimated two-year response is unstable when
that response is near zero and uncertain. Reversals of the previous
stability-only ranking cannot be dismissed as an uninformative case while
claiming a general ordering. A bounded mathematical score does not establish
welfare meaning, comparability, or reliable policy rankings. The appendix
retains the elementary boundedness result solely to document the mathematical
scope of the original construction. No GDP-loss totals, burden reductions,
or empirical trap thresholds are inferred from it.
